\bmtchapitre{Boucles et organigrammes}{Exécuter et construire une boucle, produire un organigramme et justifier la terminaison d’un algorithme simple.}{Algèbre}
\label{t11s-c07}
\begin{revision}Affectation, tests, reste de division euclidienne et calcul de sommes simples.\end{revision}
\begin{automatismes}\exo\label{t11s-c07-auto}Calculer le reste de $17$ par $5$. Si $a=3$, puis $a\leftarrow a+2$, quelle est sa valeur ? Calculer $1+2+3+4$.\end{automatismes}
\begin{corrige}[exercice~\ref{t11s-c07-auto}]\label{sol-t11s-c07-auto}Reste $2$; nouvelle valeur $5$; somme $10$.\end{corrige}

\section{Exécuter sans sauter une instruction}
\begin{definition}Une affectation $a\leftarrow E$ remplace la valeur de $a$ par la valeur actuelle de $E$. Elle ne constitue pas une égalité à résoudre. Une boucle « pour » parcourt une liste d’indices; une boucle « tant que » répète tant que sa condition est vraie. Le test de « tant que » précède chaque tour, y compris le premier.\end{definition}
\begin{exemple}[additionner les premiers entiers]
\begin{tabular}{l}
Lire un entier $n\ge0$\\
$S\leftarrow0$\\
Pour $k$ allant de $1$ à $n$\\
\quad $S\leftarrow S+k$\\
Fin pour\\
Afficher $S$
\end{tabular}
Pour $n=4$, les valeurs successives de $S$ sont $0,1,3,6,10$. Si $n=0$, aucun tour n’est effectué et la somme vide vaut $0$.
\end{exemple}
\begin{propriete}[invariant de la somme]Après le tour d’indice $k$, $S=1+\cdots+k$.\end{propriete}
\begin{demonstration}Avant le premier tour, $S=0$. Si la somme est correcte au tour précédent, l’ajout de $k$ donne celle des $k$ premiers entiers. Une récurrence justifie chaque tour; après $n$ tours, le résultat annoncé est obtenu.\end{demonstration}
\section{Un organigramme avec retour}
L’organigramme suivant calcule la même somme. Le losange est un test; les rectangles sont des affectations; les flèches indiquent l’ordre d’exécution.
\begin{visualisation}[le test décide de poursuivre]
\begin{center}\begin{tikzpicture}[node distance=9mm, every node/.style={font=\small}, process/.style={draw,rectangle,align=center,minimum width=37mm,minimum height=8mm},test/.style={draw,diamond,aspect=2,inner sep=2pt}]
\node[process](init){Lire $n\in\N$\\$S\leftarrow0$, $k\leftarrow1$};
\node[test,below=of init](test){$k\le n$ ?};
\node[process,below=of test](add){$S\leftarrow S+k$\\$k\leftarrow k+1$};
\node[process,right=20mm of test](end){Afficher $S$};
\draw[->](init)--(test);\draw[->](test)--node[right]{oui}(add);
\draw[->](test)--node[above]{non}(end);
\draw[->](add.west)--++(-12mm,0)|-(test.west);
\end{tikzpicture}\end{center}
Le compteur augmente d’une unité et finit par dépasser $n$. Le nombre de tours est fini.
\end{visualisation}
\section{Boucle conditionnelle et terminaison}
\begin{exemple}[atteindre un seuil]
Avec $u\leftarrow1$, $n\leftarrow0$, tant que $u<20$, faire $u\leftarrow2u$ puis $n\leftarrow n+1$. Les valeurs de $u$ sont $1,2,4,8,16,32$; la sortie est $n=5$. Après $n$ tours, $u=2^n$. La suite des puissances dépasse $20$ : l’arrêt est justifié.
\end{exemple}
\begin{avertissement}Si le corps ne modifie pas ce qui intervient dans le test, la boucle peut être infinie. Pour permuter deux valeurs, $a\leftarrow b$ puis $b\leftarrow a$ ne convient pas : la première valeur de $a$ a été perdue. Utiliser une variable temporaire.\end{avertissement}
\begin{methode}{contrôler un algorithme}Préciser les entrées autorisées; dresser une table d’exécution; expliquer ce qui reste vrai après chaque tour (invariant); expliquer pourquoi le test finit par être faux (terminaison); interpréter la sortie et les cas limites.\end{methode}

% ENRICHISSEMENT C07

\subsection*{Rédiger une table, une preuve et un test}
\begin{exemple}[chaque ligne utilise les valeurs actuelles]
On initialise $a=7,b=3$, puis on exécute $a\leftarrow a-b$ et $b\leftarrow a+b$.
\begin{center}\begin{tabular}{l|cc}\toprule
État&$a$&$b$\\\midrule
Initial&7&3\\Après la première affectation&4&3\\Après la seconde&4&7\\\bottomrule
\end{tabular}\end{center}
Dans la seconde ligne de calcul, $a$ vaut déjà $4$. Remplacer les deux variables simultanément par leurs anciennes valeurs décrirait un autre algorithme.
\end{exemple}
\begin{methode}{séparer résultat correct et arrêt}
Un invariant explique ce que contiennent les variables si l’algorithme atteint une étape. Il ne garantit pas que l’algorithme s’arrête. Pour une boucle bornée, compter les tours; pour une boucle conditionnelle, trouver une quantité qui finit par franchir le seuil. Tester aussi le cas où la condition est fausse dès le départ.
\end{methode}
\begin{exemple}[dénombrer les termes déjà ajoutés]
Pour calculer $u_0+\cdots+u_N$ quand $u_{n+1}=2u_n+1$, initialiser $u=u_0,S=u_0$. Au tour $k=1,\ldots,N$, mettre d’abord $u\leftarrow2u+1$, puis $S\leftarrow S+u$. Après le tour $k$, $u=u_k$ et $S=u_0+\cdots+u_k$. En $N=0$, il n’y a pas de tour et la sortie est déjà correcte. Inverser les deux mises à jour compterait deux fois $u_0$ et oublierait le dernier terme.
\end{exemple}
\begin{avertissement}[atteindre une valeur ou s’en approcher]
Une suite peut tendre vers un seuil sans l’atteindre à un rang fini. Un test d’égalité exacte ne doit pas être remplacé tacitement par un test approché. Pour demander une précision, fixer une tolérance strictement positive dans l’énoncé.
\end{avertissement}

\section{Exercices et problèmes}
\exo[Suivre les valeurs]\label{t11s-c07-e01}
Exécuter pour $n=3$ : $u\leftarrow2$; pour $k=1$ à $n$, $u\leftarrow3u-1$. Donner les valeurs et la sortie.
\begin{corrige}[exercice~\ref{t11s-c07-e01}]\label{sol-t11s-c07-e01}
Valeurs $2,5,14,41$. La sortie est $41$ après exactement trois mises à jour.
\end{corrige}
\exo[Somme de carrés]\label{t11s-c07-e02}
Écrire un algorithme calculant $1^2+\cdots+n^2$ pour $n\in\N$. Tester $n=0$ et $n=3$.
\begin{corrige}[exercice~\ref{t11s-c07-e02}]\label{sol-t11s-c07-e02}
Initialiser $S=0$; pour $k=1$ à $n$, affecter $S\leftarrow S+k^2$; afficher $S$. Sorties $0$ et $14$. Invariant : somme des carrés des indices déjà parcourus.
\end{corrige}
\exo[Échanger]\label{t11s-c07-e03}
Avec $a=4,b=7$, exécuter $a\leftarrow b$, $b\leftarrow a$. Réparer.
\begin{corrige}[exercice~\ref{t11s-c07-e03}]\label{sol-t11s-c07-e03}
On obtient $a=7,b=7$. Réparation : $t\leftarrow a$, $a\leftarrow b$, $b\leftarrow t$; résultat $a=7,b=4$.
\end{corrige}
\exo[Premier seuil]\label{t11s-c07-e04}
Exécuter $u\leftarrow3,n\leftarrow0$; tant que $u\le100$, $u\leftarrow2u$, $n\leftarrow n+1$. Quelle différence ferait le test $u<96$ ?
\begin{corrige}[exercice~\ref{t11s-c07-e04}]\label{sol-t11s-c07-e04}
Valeurs $3,6,12,24,48,96,192$; sortie $n=6$. Avec $u<96$, arrêt à $96$ et $n=5$. Le signe strict change le traitement du seuil atteint.
\end{corrige}
\exo[Une boucle sans fin]\label{t11s-c07-e05}
On initialise $u=1$ et on répète $u\leftarrow(u+1)/2$ tant que $u<2$. La boucle s’arrête-t-elle en calcul exact ?
\begin{corrige}[exercice~\ref{t11s-c07-e05}]\label{sol-t11s-c07-e05}
Non : $u=1$ reste invariant, car $(1+1)/2=1$. Le test reste vrai. La table d’un seul tour suggère l’erreur et l’invariant la prouve.
\end{corrige}
\exo[Dessiner]\label{t11s-c07-e06}
Transformer l’algorithme de la somme des carrés en organigramme avec un test avant la mise à jour. Indiquer les deux branches et l’actualisation de $k$.
\begin{corrige}[exercice~\ref{t11s-c07-e06}]\label{sol-t11s-c07-e06}
Lire $n$; initialiser $S=0,k=1$; tester $k\le n$. Oui : $S\leftarrow S+k^2$, $k\leftarrow k+1$, retour au test. Non : afficher $S$ et terminer. Les tests $n=0,3$ donnent $0,14$.
\end{corrige}

\exo[Une somme avec invariant]\label{t11s-c07-p01}
Pour un entier d’entrée $n\ge0$, on initialise $S=0,k=0$. Tant que $k<n$, on effectue, dans cet ordre, $k\leftarrow k+1$ puis $S\leftarrow S+2k$; on affiche $S$.
\begin{enumerate}
\item Dresser la table pour $n=3$ et traiter séparément $n=0$.
\item Prouver qu’après chaque tour $S=2+4+\cdots+2k$; justifier l’arrêt et interpréter la sortie.
\item Que calculerait-on si l’on ajoutait $2k$ avant d’augmenter $k$ ? Tester $n=3$.
\end{enumerate}
\begin{corrige}[exercice~\ref{t11s-c07-p01}]\label{sol-t11s-c07-p01}
Les états $(k;S)$ sont $(0;0),(1;2),(2;6),(3;12)$. Pour $n=0$, aucun tour, sortie $0$. Initialement la somme est vide; chaque tour augmente $k$ puis ajoute exactement le nouvel entier pair $2k$, ce qui conserve la description. Après $n$ tours, le test est faux et la sortie est la somme des $n$ premiers entiers pairs positifs. Avec l’ordre inversé, on ajoute $0,2,\ldots,2(n-1)$, donc la sortie pour $n=3$ vaut $6$, pas $12$.
\end{corrige}
\exo[La limite ne suffit pas à arrêter]\label{t11s-c07-p02}
En calcul exact, initialiser $u=1,n=0$ et répéter $u\leftarrow u/2$, $n\leftarrow n+1$ tant que $u>0$. Cette boucle termine-t-elle ? Modifier seulement le seuil pour demander $u\le0{,}01$, puis calculer le premier rang d’arrêt et justifier sa minimalité.
\begin{corrige}[exercice~\ref{t11s-c07-p02}]\label{sol-t11s-c07-p02}
Après $n$ tours, $u=2^{-n}>0$; aucun rang fini n’arrête la première boucle. Pour la seconde, tester tant que $u>0{,}01$. L’arrêt est atteint pour $n=7$, car $1/64>1/100$ et $1/128<1/100$. Les termes décroissent strictement; aucun rang plus petit que $7$ ne convient. La réponse concerne le calcul exact, sans hypothèse cachée sur les arrondis d’une machine.
\end{corrige}

\begin{jemeteste}Je vérifie l’initialisation, le premier et le dernier tour, le compteur et le sens du test.\end{jemeteste}
\begin{notepourlaclasse}Faire d’abord exécuter à la main. Le pseudo-code suffit; les invariants sont des justifications accessibles, sans cours de complexité. Relier ensuite aux algorithmes d’Euclide et de suites.\end{notepourlaclasse}
